\section{Risk-Trust Dynamics}
\label{sec:risk_trust}

The companion paper \citep{lasser2026horizon} introduces a risk-trust factor
$\Trisk_j$ for weighting communicated risk claims but provides only a
qualitative update rule. This section formalizes the dynamics, drawing on the
EWMA structure of the foundational paper's behavioral trust
\citep{lasser2026gfm}.

\subsection{Structural Verification Residuals}

\begin{definition}[Structural Verification Residual]
\label{def:risk_residual}
For a risk claim $(S, \Pathway, |\Delta \volL|, p)$ from agent $a_j$
    (Definition~4 of \citet{lasser2026horizon}), the actor evaluates
    \emph{the specific claimed pathway} $\Pathway$ against its own SCM
    $\SCM$ (Definition~\ref{def:scm}) using the causal risk evaluation
    procedure (Section~4.3 of \citet{lasser2026horizon}). This produces
    a pathway-conditional assessment
    $\bigl(p_{\WorldModel}(\Pathway),\;
    |\Delta \volL_{\WorldModel}(\Pathway)|\bigr)$: the probability and
    contraction magnitude the actor would assign if $\Pathway$ were
    the realized causal pathway. If $\Pathway$ has no counterpart in
    $\SCM$ --- that is, if the actor's causal model contains no edges
    implementing the claimed mechanism --- then
    $\bigl(p_{\WorldModel}(\Pathway), |\Delta \volL_{\WorldModel}(\Pathway)|\bigr)
    = (0, 0)$ by convention. The \emph{structural verification
    residual} is the dimensionless two-component divergence:
\begin{equation}
\label{eq:risk_residual}
r_j^{\mathrm{risk}}(t) = \left(
    p_j^{\mathrm{claimed}} - p_{\WorldModel}(\Pathway),\;\;
    \frac{|\Delta \volL_j^{\mathrm{claimed}}| -
          |\Delta \volL_{\WorldModel}(\Pathway)|}{\volL(G_t)}
\right)
\end{equation}
Each component is in $[-1, 1]$: the first is a probability difference,
    the second is a relative volume difference normalized by the
    current total poset volume. The Euclidean norm
    $\|r_j^{\mathrm{risk}}(t)\|$ is therefore dimensionally consistent
    and invariant under rescaling of $\volL$. A fabricated pathway
    that does not exist in $\SCM$ produces the maximal residual
    $\|r_j^{\mathrm{risk}}(t)\| = \sqrt{(p_j^{\mathrm{claimed}})^2
    + (|\Delta \volL_j^{\mathrm{claimed}}|/\volL(G_t))^2}$; structural
    verification is enforced by the pathway-conditional evaluation,
    not by an extra pathway-comparison term.
\end{definition}

Unlike the behavioral prediction residual $r_k(t)$ (Equation~25 of
\citet{lasser2026gfm}), which compares \emph{predicted} against
\emph{observed} behavior, the risk residual compares \emph{claimed} against
\emph{independently modeled} risk. This is necessary because risk claims
are counterfactual predictions that may never be directly observed —
the \emph{Wamura problem}, introduced in the poset paper
\citep{lasser2026poset} and developed in the risk-communication analysis
of \citet[\S4]{lasser2026horizon}: a successful risk
intervention removes the very evidence that would have justified it,
so the convergence timescale for risk-trust can in principle exceed
the timescale on which trust normally accumulates. The SCM provides
the independent assessment needed to break this circularity: the
actor's own causal model produces $(p_{\WorldModel},
|\Delta \volL_{\WorldModel}|)$ without relying on the claim.

\subsection{EWMA Risk-Trust Update}

\begin{definition}[Risk-Trust Dynamics]
\label{def:risk_trust_dynamics}
The risk-trust factor $\Trisk_j$ evolves through an exponentially weighted
    moving average of squared structural verification residuals:
\begin{align}
\label{eq:risk_sigma}
\sigma_j^{\mathrm{risk},2}(t) &= \alpha_{\mathrm{risk}} \cdot
    \sigma_j^{\mathrm{risk},2}(t-1) + (1 - \alpha_{\mathrm{risk}}) \cdot
    \|r_j^{\mathrm{risk}}(t)\|^2 \\
\label{eq:risk_trust_update}
\Trisk_j(t) &= \frac{1}{1 + \beta_{\mathrm{risk}} \cdot
    \sigma_j^{\mathrm{risk},2}(t)}
\end{align}
where $\alpha_{\mathrm{risk}} \in (0, 1)$ controls the decay rate
    (tracking speed vs. stability), $\beta_{\mathrm{risk}} > 0$ controls the
    sensitivity (how rapidly trust degrades with inconsistency), and
    $\|\cdot\|$ is the Euclidean norm on the two-dimensional residual.
\end{definition}

The structure mirrors the foundational paper's behavioral trust exactly
(Equations~33 and~38 of \citet{lasser2026gfm}), with the behavioral
prediction residual replaced by the structural verification residual. The
parameters $\alpha_{\mathrm{risk}}$ and $\beta_{\mathrm{risk}}$ are
independent of the behavioral parameters $\alpha$ and $\beta$: an agent can
have high behavioral trust $T_k$ (predictable behavior) and low risk-trust
$\Trisk_k$ (unreliable risk assessments), or vice versa. The two trust
channels operate independently.

\paragraph{Initialization and cooling period.} New agents enter with
$\sigma_j^{\mathrm{risk},2}(0) = \sigma_0^{\mathrm{risk},2}$ (a prior
reflecting baseline uncertainty about risk-claim quality), producing initial
risk-trust $\Trisk_j(0) = 1/(1 + \beta_{\mathrm{risk}} \cdot
\sigma_0^{\mathrm{risk},2})$. The effective cooling period is
$\tau_{\mathrm{risk}} = \lceil \log \delta / \log
\alpha_{\mathrm{risk}} \rceil$ (same formula as the behavioral cooling
period, Equation~71 of \citet{lasser2026gfm}).

\begin{proposition}[Risk-Trust $L^2$ Convergence]
\label{prop:risk_convergence}
Let $X_j(t) = \|r_j^{\mathrm{risk}}(t)\|^2$ and assume $a_j$'s risk-claim
    strategy is stationary with
    $\mu_j^{\mathrm{risk}} = \mathbb{E}[X_j] < \infty$ and
    $\nu_j^{\mathrm{risk}} = \mathrm{Var}[X_j] < \infty$. Then the
    cumulative inconsistency estimator satisfies:
\begin{align}
\label{eq:risk_trust_bias}
\lim_{t \to \infty} \mathbb{E}\bigl[\sigma_j^{\mathrm{risk},2}(t)\bigr]
    &= \mu_j^{\mathrm{risk}}
    \qquad \text{(asymptotic unbiasedness)} \\
\label{eq:risk_trust_variance}
\limsup_{t \to \infty}
    \mathrm{Var}\bigl[\sigma_j^{\mathrm{risk},2}(t)\bigr]
    &\leq \frac{1 - \alpha_{\mathrm{risk}}}
              {1 + \alpha_{\mathrm{risk}}}
       \cdot C_{\mathrm{dep}} \cdot \nu_j^{\mathrm{risk}}
    \qquad \text{(bounded stationary variance)}
\end{align}
where $C_{\mathrm{dep}} = 1$ under i.i.d.\ squared residuals, and
    $C_{\mathrm{dep}}(\alpha_{\mathrm{risk}}) = 1 + 2
    \sum_{h=1}^{\infty} \alpha_{\mathrm{risk}}^h \cdot
    \mathrm{Corr}(X_j(t), X_j(t+h))$ under weak dependence with
    summable autocovariance. The $\alpha_{\mathrm{risk}}^h$ weights
    arise from the EWMA's geometric weighting: the cross-terms in the
    variance expansion carry the product of the EWMA coefficients at
    lag $h$. Since $|\mathrm{Corr}| \leq 1$ and
    $\sum_h \alpha_{\mathrm{risk}}^h = \alpha_{\mathrm{risk}} /
    (1 - \alpha_{\mathrm{risk}})$, we have the universal bound
    $C_{\mathrm{dep}} \leq (1 + \alpha_{\mathrm{risk}}) /
    (1 - \alpha_{\mathrm{risk}})$.
Consequently, the variance estimator
    $\sigma_j^{\mathrm{risk},2}(t)$ concentrates in $L^2$ around
    $\mu_j^{\mathrm{risk}}$. The risk-trust factor $\Trisk_j(t)$
    concentrates around the nominal fixed point
    $\Trisk_j^* = 1 / (1 + \beta_{\mathrm{risk}} \mu_j^{\mathrm{risk}})$
    in distribution (via the continuous mapping theorem), with the
    caveat that $\mathbb{E}[\Trisk_j(t)]$ may exceed $\Trisk_j^*$ by
    a Jensen bias term of order
    $O(\beta_{\mathrm{risk}}^2 \cdot
    \mathrm{Var}[\sigma_j^{\mathrm{risk},2}])$ due to the concavity
    of $x \mapsto 1/(1 + \beta_{\mathrm{risk}} x)$. The deviation
    is controlled by $\alpha_{\mathrm{risk}}$: as
    $\alpha_{\mathrm{risk}} \to 1$, the stationary variance vanishes,
    the Jensen bias vanishes, and $\Trisk_j(t) \to \Trisk_j^*$ in
    $L^2$. Structurally accurate reporters (small
    $\mu_j^{\mathrm{risk}}$) concentrate around high $\Trisk_j^*$;
    alarmists and chronic underreporters concentrate around low
    $\Trisk_j^*$.
\end{proposition}

\begin{proof}
Unrolling the EWMA recursion gives
$\sigma_j^{\mathrm{risk},2}(t) = (1 - \alpha_{\mathrm{risk}})
    \sum_{s=0}^{t-1} \alpha_{\mathrm{risk}}^s X_j(t-s)
    + \alpha_{\mathrm{risk}}^t \sigma_j^{\mathrm{risk},2}(0)$.
Taking expectations under stationarity,
$\mathbb{E}[\sigma_j^{\mathrm{risk},2}(t)]
    = (1 - \alpha_{\mathrm{risk}})
      \sum_{s=0}^{t-1} \alpha_{\mathrm{risk}}^s \mu_j^{\mathrm{risk}}
    + \alpha_{\mathrm{risk}}^t \sigma_j^{\mathrm{risk},2}(0)
    \to \mu_j^{\mathrm{risk}}$ as $t \to \infty$, establishing
    Equation~\eqref{eq:risk_trust_bias}. Under the i.i.d.\ assumption,
\begin{align*}
\mathrm{Var}[\sigma_j^{\mathrm{risk},2}(t)]
    &= (1 - \alpha_{\mathrm{risk}})^2
       \sum_{s=0}^{t-1} \alpha_{\mathrm{risk}}^{2s}
       \cdot \nu_j^{\mathrm{risk}}
       + o(1) \\
    &\longrightarrow \frac{(1 - \alpha_{\mathrm{risk}})^2}
                          {1 - \alpha_{\mathrm{risk}}^2}
       \cdot \nu_j^{\mathrm{risk}}
       = \frac{1 - \alpha_{\mathrm{risk}}}
              {1 + \alpha_{\mathrm{risk}}}
       \cdot \nu_j^{\mathrm{risk}},
\end{align*}
giving $C_{\mathrm{dep}} = 1$ in
Equation~\eqref{eq:risk_trust_variance}. Under weak dependence
with summable autocovariance, the cross-terms in the EWMA variance
expansion carry $\alpha_{\mathrm{risk}}^h$ weights from the
geometric kernel, giving
$C_{\mathrm{dep}}(\alpha_{\mathrm{risk}}) = 1 + 2
\sum_{h=1}^{\infty} \alpha_{\mathrm{risk}}^h \cdot
\mathrm{Corr}(X_j(t), X_j(t+h))$, which is finite because
$\alpha_{\mathrm{risk}}^h$ decays geometrically and the
autocovariance is summable. The
concentration statement for $\sigma_j^{\mathrm{risk},2}(t)$ around
$\mu_j^{\mathrm{risk}}$ in $L^2$ follows directly. For $\Trisk_j(t)$,
the continuous mapping theorem applied to
$x \mapsto 1 / (1 + \beta_{\mathrm{risk}} x)$, which is Lipschitz on
$[0, \infty)$, gives convergence in distribution to
$\Trisk_j^*$. The Jensen bias arises because
$\mathbb{E}[f(X)] \geq f(\mathbb{E}[X])$ for concave $f$; a
second-order Taylor expansion gives the bias as
$O(\beta_{\mathrm{risk}}^2 \cdot
\mathrm{Var}[\sigma_j^{\mathrm{risk},2}])$, which vanishes as
$\alpha_{\mathrm{risk}} \to 1$.
\end{proof}

\paragraph{Tracking versus convergence: why $L^2$ and not almost sure.}
Fixed-$\alpha$ EWMA does not converge almost surely to a deterministic
limit, and for an adversarial-tracking setting this is the correct
behavior. Each new residual contributes weight $(1 - \alpha_{\mathrm{risk}})$
regardless of how much history has accumulated, so the estimator
never collapses to a point mass and retains nonzero stationary
variance. A Robbins-Monro stochastic approximation with step size
$\alpha_t \to 0$ satisfying $\sum_t \alpha_t = \infty$ and
$\sum_t \alpha_t^2 < \infty$ \citep{robbins1951stochastic} would
converge almost surely to $\mu_j^{\mathrm{risk}}$, but only at the
cost of losing the ability to track a non-stationary adversary whose
strategy changes mid-trajectory. Risk-claim quality and scorpion
behavior are both moving targets: a cooperative reporter can become
systematically biased, an alarmist can calibrate, and a scorpion can
change strategy in response to detection pressure. Past accuracy does
not guarantee future accuracy under a changing world model
$\WorldModel$, so the tracking property is load-bearing. The
stationary-variance bound in Equation~\eqref{eq:risk_trust_variance}
is what the scorpion detection threshold
(Section~\ref{sec:scorpion}) is calibrated against, so the $L^2$
form is not just weaker than almost sure --- it is the form the
downstream results actually use.

\paragraph{Interaction with the Wamura problem.} The EWMA dynamics converge
on \emph{structural verification quality}, not on \emph{outcome correctness}.
An agent whose risk claims are structurally plausible (the pathway exists in
$\SCM$, the probability and magnitude estimates are calibrated against the
actor's own model) accumulates low $\|r^{\mathrm{risk}}\|^2$ and converges
to high $\Trisk_j^*$ regardless of whether the risk materializes. This
partially resolves the Wamura problem: risk-trust converges even when the
counterfactual is never observed, because convergence depends on structural
agreement, not outcome observation. The resolution is partial because an
agent whose claims are structurally plausible but whose probability estimates
are systematically biased (always 10\% higher than the actor's model)
accumulates moderate $\|r^{\mathrm{risk}}\|^2$ and converges to moderate
$\Trisk_j^*$, reflecting the bias without resolving it.

\subsection{Independence-Preserving Risk Aggregation}

The max operator in Equation~7 of \citet{lasser2026horizon} destroys
independence information: 50 low-trust agents agreeing on a risk assessment
are treated identically to one low-trust agent. We propose a
trust-weighted log-odds aggregation that preserves the evidence-accumulation
structure of independent reports.

\paragraph{Definition update from the companion paper.}
The companion paper \citep{lasser2026horizon} defines
$\hat{\Risk}_j$ as an expected-contraction product
$T_j^{\mathrm{risk}} \cdot p \cdot |\Delta \volL|$, which is not a
probability and cannot be passed through log-odds. We update the
definition: $\hat{\Risk}_j$ and $\Risk_{\mathrm{self}}$ hereafter
denote posterior probabilities that the exercise pathway will be
executed within the planning horizon, calibrated against the shared
prior $\pi_0$ below. The contraction magnitude $|\Delta \volL|$
becomes a separate quantity used by the verification gate
(Section~\ref{sec:attention_not_action}) and the structural
verification residual (Definition~\ref{def:risk_residual}). The
canonical form of a communicated risk claim is the pair
$(\hat{\Risk}_j, |\Delta \volL_j^{\mathrm{claimed}}|)$; the
single-number product in \citet{lasser2026horizon} is superseded.

\begin{definition}[Log-Odds Risk Aggregation]
\label{def:logodds}
Let $\pi_0 \in (0, 1)$ be a shared prior on the probability that the
    exercise pathway will be executed within the planning horizon, and
    let $\ell_0 = \log(\pi_0 / (1 - \pi_0))$ be its log-odds. The
    actor's self-assessment $\Risk_{\mathrm{self}}$ and each communicated
    risk estimate $\hat{\Risk}_j$ from agent $a_j$ are posterior
    probabilities derived from the shared prior $\pi_0$ plus their
    respective observations (definition update, see above). The aggregated posterior is:
\begin{equation}
\label{eq:logodds}
\log \frac{\Risk_{\mathrm{agg}}}{1 - \Risk_{\mathrm{agg}}}
    = \ell_0
    + \underbrace{\left(\log \frac{\Risk_{\mathrm{self}}}{1 - \Risk_{\mathrm{self}}} - \ell_0\right)}_{\text{actor's log-likelihood-ratio}}
    + \sum_j \Trisk_j \cdot \underbrace{\left(\log \frac{\hat{\Risk}_j}{1 - \hat{\Risk}_j} - \ell_0\right)}_{\text{reporter $a_j$'s log-likelihood-ratio}}
\end{equation}
where each parenthesized term is the log-likelihood-ratio contributed
    by that source beyond the shared prior, obtained by subtracting
    $\ell_0$ from their reported posterior log-odds. The $\Risk$ values
    are normalized to $[0, 1]$ as probabilities of the exercise pathway
    being executed within the planning horizon. The shared prior
    $\pi_0$ is an operational coordination requirement: every
    reporter's posterior must be calibrated against the same $\pi_0$
    for the subtraction $\log(\hat{\Risk}_j / (1 - \hat{\Risk}_j))
    - \ell_0$ to correctly extract the log-likelihood ratio.
    Establishing $\pi_0$ requires a protocol (e.g., broadcasting it
    as a system parameter at population initialization), not just a
    mathematical assumption.
    Under the neutral
    bootstrap prior $\pi_0 = 0.5$ ($\ell_0 = 0$), the prior-correction
    terms vanish and Equation~\eqref{eq:logodds} reduces to a naive
    sum of posterior log-odds; under any other shared prior, the
    correction is required to avoid double-counting the prior.
\end{definition}

The log-odds aggregation has three advantages over the max operator:
\begin{enumerate}
\item \textbf{Agreement strengthens evidence.} Multiple independent
    low-trust agents reporting the same risk produce a collective signal
    proportional to the number of agreeing reporters, weighted by their
    individual $\Trisk_j$ values.
\item \textbf{Trust-weighted log-linear opinion pool.}
    Equation~\eqref{eq:logodds} is a trust-tempered log-linear opinion
    pool \citep{genest1986combining}: the aggregate log-odds is a
    weighted sum of the reporters' log-likelihood ratios plus the
    shared prior. Under the conditional-independence assumption
    (reporters' observations are independent given the true risk
    level and the shared prior $\pi_0$), setting all $\Trisk_j = 1$
    recovers the standard Bayesian posterior update for independent
    evidence sources. The trust weights $\Trisk_j \in [0, 1]$ temper
    each reporter's contribution in proportion to their
    structural-verification track record; this is not itself
    Bayes-optimal in general --- exact Bayes-optimality would require
    a generative model of how each reporter's noise depends on their
    history, which the framework does not assume --- but it recovers
    Bayes-optimal behavior in the no-tempering limit and degrades
    gracefully for untrusted reporters, which is the behavior the
    aggregation rule is trying to achieve. The prior-correction
    terms are required regardless of tempering: without them, a naive
    sum of posterior log-odds would double-count $\pi_0$ once per
    reporter.
\item \textbf{Graceful degradation.} When a single agent has
    $\Trisk_j \gg \sum_{k \neq j} \Trisk_k$, the aggregation approximates
    max behavior: the dominant reporter's assessment controls the aggregate.
    The max operator emerges as a limiting case, not an imposed design
    choice.
\end{enumerate}

\paragraph{Bootstrap interaction.} During bootstrap, all $\Trisk_j$ values
start at the neutral prior. Under log-odds aggregation, $n$ agents
reporting the same risk at neutral trust produce a collective signal
$n \times \Trisk_{\mathrm{neutral}}$ times the individual log-odds,
proportional to $n$. Under the max operator, the same $n$ agents produce
only $\Trisk_{\mathrm{neutral}}$ times the maximum individual log-odds,
independent of $n$. The log-odds aggregation therefore learns faster during
bootstrap, precisely when the actor's own model is least reliable and
external evidence is most valuable.

\subsection{Aggregation as Attention, Not Action}
\label{sec:attention_not_action}

The log-odds aggregation amplifies agreement among reporters, which is
correct when reporters are independent but exploitable by coordinated
conspirators: $n$ agents who agree in advance on a fabricated risk claim
produce a collective signal proportional to $n$, potentially overwhelming
the actor's self-assessment and triggering restriction of a capability that
was never actually dangerous. This is the risk-communication analogue of
report-bombing: coordinated low-quality reports exploiting an aggregation
mechanism designed for independent evidence.

Any fixed correlation-detection heuristic is game-able once the adversary
knows the parameters --- a standard observation in adversarial detection
settings --- and the ROC trade-off is unfavorable in this domain because
both error types are costly: false positives (restricting a safe
capability) and false negatives (ignoring a real risk) both cause
$\volL$-contraction. We therefore adopt the activation
energy principle from the foundational paper \citep{lasser2026gfm}, applied
to risk-claim processing.

\paragraph{The two-gate architecture.} The risk aggregation
$\Risk_{\mathrm{agg}}$ (Definition~\ref{def:logodds}) does not directly
trigger restriction. It produces a \emph{salience signal}: the aggregated
risk level determines the priority at which a capability tuple enters the
actor's investigation queue. Action requires passing two independent gates:
\begin{enumerate}
\item \textbf{Salience gate.} The aggregated risk
    $\Risk_{\mathrm{agg}}(S; a_k, t)$ exceeds an attention threshold
    $\theta_{\mathrm{risk}}$, placing the capability tuple in the
    investigation queue.
\item \textbf{Verification gate.} The actor invokes the forward
    causal risk-evaluation procedure (Section~4.3 of
    \citet{lasser2026horizon}) on its own SCM $\SCM$ to independently
    confirm that the claimed exercise pathway $\Pathway$ exists in
    the actor's causal model, calibrate the probability, and verify
    the contraction magnitude. This is a forward causal query ---
    ``if $\Pathway$ were executed, what would happen?'' --- using
    the same pathway-conditional assessment
    $\bigl(p_{\WorldModel}(\Pathway), |\Delta \volL_{\WorldModel}(\Pathway)|\bigr)$
    as the structural verification residual
    (Definition~\ref{def:risk_residual}). It is distinct from the
    retrospective causal contraction attribution
    (Definition~\ref{def:causal_attribution}), which decomposes an
    \emph{observed} contraction across agents after the fact.
    Retrospective attribution has no counterfactual access to
    hypothetical pathways that were never executed; forward risk
    evaluation does, and this is the mechanism the verification gate
    requires. The gate is the activation energy: the actor must
    reconstruct the risk through its own causal model before
    acting.
\end{enumerate}
Neither gate alone is sufficient. A coordinated report-bomb raises salience
(gate~1) but fails independent verification (gate~2) because the fabricated
pathway does not exist in the actor's causal model. A genuine risk reported
by a single low-trust agent may fail gate~1 (low salience) but would pass
gate~2 if the actor investigated. The two-gate architecture trades a
bounded false-negative rate (genuine risks from low-trust reporters may be
investigated slowly) for strong robustness against coordinated manipulation
(fabricated risks cannot bypass independent verification regardless of how
many conspirators report them).

\paragraph{Why heuristic correlation detection is insufficient.} A fixed
correlation metric (temporal clustering, behavioral similarity among
reporters, trust-history patterns) is game-able by a sufficiently
informed adversary: adversaries who know the detection threshold can
space their reports, diversify their trust-building histories, and
avoid clustering. The two-gate architecture avoids this arms
race by not relying on detecting coordination at all. The defense is the
actor's own intelligence applied through the causal model, not a statistical
filter on the report stream.
